\section{Prólogo a la primera edición}\label{prologo-a-la-primera-edicion}

Aunque bien lo advertirá el lector a poco que hojee este volumen, no estará de más afirmar desde luego que se trata de un \emph{Manual de historia de España}, es decir, de un libro elemental, de vulgarización, que no tiene pretensiones eruditas, ni presume de agotar la materia, ni mucho menos de enseñar nada a los estudiosos, familiarizados ya con todas y cada una de las relativas novedades que para cierta parte del público seguramente contiene. Al escribirlo, se ha pensado ante todo en ese público, falto de tiempo y de preparación para leer obras extensas o de carácter crítico, como para enfrascarse en la ardua tarea de estudiar monografías e ir traduciendo luego, poco a poco, el conjunto de los resultados parciales, en conclusiones de alcance general; y también se han tenido en cuenta las necesidades de una gran masa escolar que cada día exige con mayor imperio libros acomodados a los modernos principios de la historiografía y a los progresos indudables que la investigación ha realizado, de pocos años a esta parte, en lo que se refiere a la vida pasada del pueblo español.

No quiero decir con ello que la literatura histórica de nuestra patria carezca de libros de este género, a tal punto que pueda ofrecerse el actual como novedad sin precedentes. Comienzo, por el contrario, afirmando que soy un mero continuador de ensayos anteriores valiosos, un obrero más que intenta, a su modo y con las pobres fuerzas de que dispone, resolver nuevamente el problema de un \emph{Manual de historia de España} que pueda servir para la enseñanza en varios de sus grados y para la cultura general, necesitada aquí, como en ninguna otra parte, de libros de escaso volumen, de fácil lectura, de poco aparato científico y de moderado precio, y que, juntamente, se amolden a los principios metodológicos seguidos hoy día en todos los países, conforme el propio autor ha expuesto en otro lugar\footnote{Rafael Altamira y Crevea, \emph{La enseñanza de la historia}, 2.ª ed. (Madrid, 1895).}.

En consideración a esos principios, de gloriosa tradición nacional, se ha titulado el libro \textsc{Historia de España y de la civilización española}, para evitar que, llamándose a secas \emph{Historia de España}, se creyese que solo comprendía (como es uso corriente) la parte política externa, o que, adoptando tan solo el nombre de \emph{Historia de la civilización española}, excluía ---como muchas obras que se apellidan así--- aquella parte tan esencial en la vida de los pueblos, reduciéndose a pura historia interna del movimiento civilizador que, además, no todos los autores entienden de igual modo. Continuando la difusión de las ideas (que podemos llamar \emph{modernas} no obstante su antiguo abolengo, puesto que solo en nuestros días han adquirido aceptación universal y se han formulado sistemáticamente) acerca del concepto y el contenido de la historia, llegará momento en que baste decir historia de tal o cual nación para que se entienda por todos que comprende, tanto las manifestaciones externas como las internas de la actividad social. Hoy por hoy, aún me parece oportuno dirigir la atención del lector con esos apelativos mixtos, que ya usó nuestro gran Masdeu; porque, no obstante la inclusión en obras extensas, como la de Lafuente\footnote{La \emph{Historia general de España} de Modesto Lafuente, disponible en la \href{https://catalogo.bne.es/permalink/34BNE_INST/95vni4/alma991045007239708606}{Biblioteca Nacional de España}, se publicó entre 1850 y 1867 en 30 volúmenes, y es obra de referencia ineludible en la erudición española del siglo \textsc{xix}. Altamira la cita como ejemplo paradigmático de una historia que, pese a incluir capítulos sobre la civilización, mantiene la política como eje narrativo central.}, de capítulos relativos a la civilización, por ser estos de mucho menor desarrollo que los dedicados a la historia política externa y sin la debida proporción con ellos, la mayoría de los lectores sigue entendiendo a la manera antigua el contenido de la narración histórica.
